\chapter{Introduction}
\label{ch:intro}

\section{Motivation}

Graphs describe road and transport networks, power grids, dependency structures, social and financial networks.
Many everyday questions about such systems are graph questions: is there still a route if this station closes, what
is the shortest way from A to B, does this set of tasks contain a circular dependency. Large language models (LLMs)
are increasingly asked questions like these in natural language, and a body of work has studied how well they answer
them.

The standard way of giving an LLM a graph is to serialize it: the edge list or adjacency list is written into the
prompt, and the model is asked a question about it. The prior work this thesis builds on, Skianis et al.\
\cite{skianis2024pseudocode}, followed this route and showed that the way the model is prompted matters: adding
pseudo-code of the relevant algorithm, alone or with a worked example, changes how well models reason over the
serialized graph. Serialization has two structural problems, however. First, the prompt grows with the graph, so the
cost of every question grows with it and large graphs do not fit at all. Second, the model is left to do work it is
poor at: counting long lists, searching exhaustively, and copying dozens of node identifiers without a mistake.

Agent harnesses offer a different arrangement. A harness is the program around a model: it decides what the model
sees, which tools it may call, when a run stops, and what happens to the answer. General-purpose harnesses such as
Claude Code \cite{claudecode} let a strong model write and run code, so a graph question can be answered by a
networkx \cite{networkx} script instead of by reading the graph. This works, but it is expensive: the harness carries
a large system prompt and tool machinery built for software engineering, and it relies on a frontier model that
knows how to use a graph library. In a small preview run (Chapter~\ref{ch:results}), Claude Code used on average about
41,900 input tokens per question on graph questions that our harness answered with about 1,500.

This thesis asks whether a harness built specifically for graph problems can get the same answers from a cheap
model. The graph stays in Python memory and is never written into the prompt; the model sees only the results of
validated graph tools. The harness checks answers with code, keeps every message bounded regardless of graph size,
and logs everything, so that accuracy and cost can be measured and reproduced.

\section{Research question}

The core research question is:

\begin{quote}
Can a graph-specific harness let a cheap model match or beat a frontier model running in a general-purpose harness
(e.g.\ Claude Code), at a fraction of the cost?
\end{quote}

``Better'' is judged on the accuracy-versus-cost Pareto front, not on raw accuracy alone: a configuration that is
slightly less accurate but an order of magnitude cheaper can be the better choice, and a configuration is only
dominated if another one is at least as accurate and at least as cheap. Cost is measured in tokens and latency for the
local models, where the dollar cost is zero, and in dollars (API-equivalent) for hosted models.

The question breaks down into sub-questions that follow the components of the harness:
\begin{enumerate}
  \item How accurate is a cheap local model when it solves graph problems through tools, compared with the same model
        and with a frontier model in a general-purpose harness given the same tools?
  \item What does each harness component contribute (tool design, verifiers, result handles, code execution,
        routing and skills, escalation) to accuracy and cost? This is answered by ablations.
  \item Where do cheap models fail inside a harness, and which failures can the harness remove? This is a
        process-level question (the proposal's Option B), answered from the trajectories rather than the final answers.
\end{enumerate}
\todo[inline]{Check the sub-questions against the proposal (docs/thesis\_proposal\_options\_ABC.pdf). The PDF is referenced in CLAUDE.md but is not in the repository on 2026-10-09, so this draft could not use it.}

\section{Contributions}

The thesis is in progress. The contributions below are separated into what is built and measured on the date of this
draft and what is planned.

\paragraph{Achieved so far.}
\begin{itemize}
  \item A working graph-specific agent harness in Python (Chapters~\ref{ch:design} and~\ref{ch:impl}): an agent loop
        with budgets, loop detection and rescue of malformed tool calls; 13 default graph tools and 5 opt-in tools,
        each validating its input and never raising; per-task structured answers; result handles for large outputs;
        a bound on every message the model can see; context compaction; prompt versioning; and full JSONL tracing.
  \item Evidence-based verifiers: programmatic checks that inspect the evidence an answer carries (a path, a cycle, a
        spanning forest, a neighbor list) and do not re-solve the question, together with an explicit account of which
        answers can and cannot be checked this way.
  \item An MCP server that exposes exactly the same tools, so that general-purpose harnesses can be compared with
        ours on equal tools, and a locked-down Docker sandbox for model-written code.
  \item An evaluation toolkit: a runner over the existing dataset and generated tasks, a report with strict and
        lenient accuracy and bootstrap confidence intervals, and process metrics (tool precision, recall, F1,
        parameter match) defined as in the GDS Agent benchmark \cite{gdsagent,gdsagentbench}.
  \item Preliminary findings (Chapter~\ref{ch:results}): on the existing dataset a cheap model with fitting tools is
        accurate, and its remaining errors are formatting errors rather than reasoning errors; well-fitting tools cut
        cost by roughly an order of magnitude at equal accuracy; on aggregation tasks, code over the same tools
        succeeds where tool-by-tool calling fails; and a silent serving setting (a 4,096-token context) can masquerade
        as model behaviour.
\end{itemize}

\paragraph{Planned.}
\begin{itemize}
  \item A router and per-family skills, with per-family tool filtering (M4), and an interactive command-line
        interface (M4.5).
  \item Escalation to a stronger model when verification fails (M5).
  \item Benchmark adapters (GABench \cite{gabench}, ProGraph \cite{prograph}, GrAlgoBench \cite{gralgobench},
        GT Bench \cite{gtbench}, and the GDS Agent benchmark \cite{gdsagentbench}) and baselines in the same
        conditions: pure prompting, a generic ReAct-style agent with Python, Claude Code headless, GraphTeam
        \cite{graphteam} (M6).
  \item The final experiments: size scaling, Pareto fronts, component ablations, and repeated runs with confidence
        intervals (M7).
\end{itemize}

\section{Scope and status of this draft}

All results in this draft come from the development split of the dataset, most from a single run, on small question
sets, with confidence intervals only where the report computes them. The project plan requires at least three runs
per configuration and bootstrap confidence intervals before any result is treated as final; this draft does not meet
that bar yet. The results are reported to document the state of the work and the lessons learned while building the
harness, not as the thesis's answer to the research question.

\section{Structure}

Chapter~\ref{ch:background} reviews LLM-based graph reasoning, tool use and agent harnesses, and graph benchmarks.
Chapter~\ref{ch:design} presents the design principles and the architecture. Chapter~\ref{ch:impl} describes the
implementation as built. Chapter~\ref{ch:setup} describes the dataset, models, configurations and metrics.
Chapter~\ref{ch:results} reports the results so far. Chapter~\ref{ch:discussion} discusses them and collects the
lessons learned. Chapter~\ref{ch:remaining} lists the remaining work, and Chapter~\ref{ch:conclusion} gives a
provisional conclusion. The appendices list the tools and the text the harness sends to the model.
